\documentclass{report}
\usepackage{amsmath,amssymb}
\setlength{\parindent}{0mm}
\setlength{\parskip}{1em}
\begin{document}
\begin{center}
\rule{6in}{1pt} \
{\large Alyson, L Fox \\
{\bf Numerical Methods for Gremban's Expansion of Signed Graphs}}

782 30th Street \\ Boulder \\ CO 80303
\\
{\tt fox.alyson@gmail.com}\end{center}

Signed graphs are an interesting subset of the typical graph problem.
They not only show how two vertices can be positively related but also
how negative relations as well. Data scientists have studied linear
solvers for unsigned, undirected graph Laplacians at length, but the
supplementary knowledge of signed graphs, gives current solvers
difficulty. Gremban's Expansion \cite{GrembanPHD} can be used to
transform the signed, undirected graph Laplacian into an unsigned,
undirected graph Laplacian. The solution to the linear system of the
expanded matrix yields the solution to the linear system of the original
matrix. Unsigned, undirected graph Laplacian solvers have been shown to
be robust, thus, using Gremban's Expansion, current solvers can also
solve the signed, undirected graph Laplacian linear system. This paper
delves into the numerical stability and applicability of Gremban's
expansion for signed graph Laplacians. A tight error bound for Gremban's
expansion for a symmetric linear system is shown. It is also shown that
Gremban's expansion can be extended to nonsymmetric graph Laplacians.
Both manufactured and real-world signed, undirected graph Laplacians are
tested with various solvers to show that the expansion is numerically
stable.


\end{document}
